\section{General conventions and probabilistic tools}
\label{sec:general-conventions-and-probabilistic-tools}

This section supplies the finite tools used throughout the proof. We first
retain witness families after small removals and balance their expected
label loads. We then prove conditional probability comparisons, a local
construction for choosing centers, and two samplers that assign distinct
labels. Each later application specifies the laws and local tests to which
these tools are applied.

\begin{definition}[Availability]
A witness has specified supports on the current host sides. After a pair
of label removals, it remains a witness if both supports are wholly
retained. A family is \textbf{available} along a sequence if some constant
\(\kappa>0\) has the following property: eventually, after every removal
of at most \(\kappa N\) labels from each side, at least one witness remains.
The witness may be chosen anew for each pair of removals.
\end{definition}

\begin{lemma}[Stabilization and finite menus]\label{lem:stabilization}
For a countable list of existence properties defined by witnesses
intrinsic to their supports, we can pass to a subsequence and discard
\(o(N)\) labels so that each property is either eventually absent or
available. Every fixed earlier availability survives with a smaller fixed
removal tolerance. An available finite union has an available member after
passing to a subsequence. An available menu can be replaced by a finite
menu at each dimension.
\end{lemma}

\begin{proof}
Treat the properties in order. If the current property can be eliminated
by \(o(N)\) additional removals along a subsequence, make those removals
and pass to that subsequence. Otherwise it is available: failure of every
fixed tolerance would let us choose successive dimensions at which
removals of size at most \(N/j\) eliminate the property, producing just
such a subsequence. Diagonalize these nested subsequences, choosing each
term sufficiently far out that the accumulated removals have size
\(o(N)\). For each fixed earlier stage, the later removals eventually use
less than half of its tolerance. Its availability therefore persists with
half that tolerance, while an absence persists under every further removal.

For an available union of \(k\) families with tolerance \(\kappa\), use
\(\kappa/k\) as the removal tolerance for each member. If no member were
available on a subsequence, we could combine their eliminating sets and
eliminate the union within its tolerance. Finally, there are only finitely
many pairs of label subsets at any fixed dimension. For each allowed
removal pair choose one surviving witness. These representatives form a
finite menu with the same availability.
\end{proof}

\begin{lemma}[Balanced mixtures and simultaneous profiles]\label{lem:balanced-profiles}
Let a finite patch menu be available with tolerance \(\kappa>0\), and let
patch \(i\) carry laws \((\mu_i,\nu_i)\) on its two supports. There is a
probability distribution on the menu such that, coordinatewise,
\[
 \mathbb E_i\mu_i,\ \mathbb E_i\nu_i\le K/N
\]
for a constant \(K\).

More generally, let finitely many players have finite action sets and
choose their actions independently according to probability distributions
called \emph{profiles}. Conditional on the realized actions, fix each
player's output-vector rule independently of the profiles. Suppose that,
for every choice of the other players' profiles, a player can meet its
fixed expected-vector upper bounds by a mixture of its actions. Then
there is a profile for every player under which all these bounds hold
simultaneously.
\end{lemma}

\begin{proof}
For the first assertion, assign nonnegative prices to labels on the
disjoint union of the host sides, and let \(P\) be their total. If \(P=0\)
there is nothing to prove. Otherwise, fewer than \(\kappa N\) labels on
either side have price greater than \(2P/(\kappa N)\). Availability gives
a patch avoiding all these labels. Each of its two laws then has expected
price at most \(2P/(\kappa N)\), so its pair-vector has price at most
\(4P/(\kappa N)\). If the convex hull of the pair-vectors contained no
vector bounded coordinatewise by \(4/(\kappa N)\), separation from that
downward-closed region would give nonnegative prices contradicting this
inequality. Thus \(K=4/\kappa\) suffices.

For the second assertion, write \(X_j\) for player \(j\)'s output vector
and \(b_j\) for its prescribed upper bound. With the other profiles
\(p_{-j}\) fixed, its feasible responses form
\[
 \mathcal R_j(p_{-j})
   =\{q_j:\mathbb E_{q_j,p_{-j}}X_j\le b_j\}.
\]
These sets are nonempty by hypothesis, and are closed and convex because
expectation is linear in \(q_j\). Since the action sets are finite and the
output rules are fixed, the expectations are continuous in all profiles;
hence the response correspondence has closed graph. The product of these
correspondences has a fixed point by Kakutani's theorem~\cite{kakutani1941}.
At that point every player's expected-vector bound holds.
\end{proof}

A kernel below may be a \emph{subprobability}, a nonnegative measure of
mass at most one. For example, a row \(p_v\), indexed by a cube role
\(v\), may be set to zero when its local construction is invalid.
Auxiliary experiments may instead use a specified fallback draw to
continue sampling. The final rows on a successful outcome must be
probability measures. If they are supported on the allowed labels and
\(\sum_v p_v(y)\le1\) for each label \(y\), then for every set \(S\) of
roles,
\[
 |S|=\sum_y\sum_{v\in S}p_v(y)
 \le\left|\bigcup_{v\in S}\operatorname{supp}p_v\right|.
\]
Hall's theorem~\cite{hall1935} therefore gives distinct allowed
representatives.

\begin{definition}[Raw experiments, gates, and stage histories]
A \emph{raw experiment} is the specified sampling law before the avoidance
constraints currently under discussion are imposed. A \emph{local gate}
is an event determined by the declared local inputs; retaining the gate
means multiplying by its indicator, rather than conditioning the
reference law on it. A \emph{successful prehistory} is an entering stage
history on which that stage's stated hypotheses hold. Comparisons are
made separately at each such history.
\end{definition}

The next comparisons bound expectations of arbitrary nonnegative tests.
An absolute estimate for the probability of failure would not suffice
when the test itself has very small probability.

\begin{lemma}[Conditional avoidance comparison]\label{lem:conditional-avoidance}
Let \(E_i\) be finitely many bad events with a symmetric neighbor relation
\(\sim\). Suppose that, whenever defined, the probability of \(E_i\)
conditioned on avoidance of any subset of its nonneighbors other than
itself is at most \(p_i\). Independent underlying variables, with events
joined when their variable scopes overlap, satisfy this hypothesis with
\(p_i=\Pr(E_i)\). Assume
\[
 p_i\le x_i\prod_{j\sim i}(1-x_j),\qquad 0\le x_i<1.
\]
Then all the events can be avoided with positive probability. For any set
\(S\) of indices not containing \(i\),
\(\Pr(E_i\mid \bigcap_{j\in S}E_j^c)\le x_i\).

Write \(A_S=\bigcap_{j\in S}E_j^c\). For disjoint index sets \(S,T\) and
any nonnegative random variable \(W\),
\[
 \mathbb E[W\mid A_{S\cup T}]
 \le \prod_{j\in T}(1-x_j)^{-1}\mathbb E[W\mid A_S].
\]
If an additional event \(F\) has conditional probability at most \(p_F\)
under avoidance of any subset of its nonneighbors, then
\[
 \Pr(F\mid A_S)\le
 p_F\prod_{j\in S:\,j\sim F}(1-x_j)^{-1}.
\]
\end{lemma}

\begin{proof}
The local lemma and its conditional induction go back to Erd\H{o}s and
Lov\'asz~\cite[Section~2]{erdos-lovasz1975}. Erd\H{o}s and
Spencer introduced the lopsided nonneighbor condition~\cite[Section~1]{erdos-spencer1991};
see Lu and Sz\'ekely~\cite[Lemma~3]{lu-szekely2007} for the asymmetric
formulation used here.
Here is the induction for the stated hypotheses. Split a conditioning
set into nonneighbors and neighbors of \(E_i\). The numerator is bounded
by \(p_i\) after discarding neighbor-avoidance conditions. By induction,
each neighbor-avoidance condition has conditional probability at least
\(1-x_j\), so the denominator is at least their product. This proves
\(\Pr(E_i\mid A_S)\le x_i\), and simultaneously proves positivity of all
avoidance probabilities. In particular,
\[
 \Pr(A_T\mid A_S)\ge\prod_{j\in T}(1-x_j).
\]
Inserting \(\mathbf1_{A_T}\le1\) in the conditional expectation of
\(W\), then dividing by this lower bound, proves the first comparison.
The same argument for \(F\) removes only its neighbors and uses its own
conditional nonneighbor bound for the remaining numerator.
\end{proof}

For independent underlying variables, removing all avoidance constraints
that touch a specified set of variables makes those variables independent
of the remaining avoidance event. They can then be integrated under their
raw laws. In applications we impose avoidance separately at each entering
history; we do not reweight earlier draws by the probability of later
avoidance.

\begin{lemma}[Scattered moments]\label{lem:scattered-moments}
Let \(Z_v\ge0\), indexed by a finite role set \(U\), be bounded by \(L\)
on a success event. Each role has a near set, including itself, of size at
most \(f|U|\). Suppose that for any sequence of at most \(n\) indices,
each outside the near sets of its predecessors, the joint product
expectation with the success indicator is at most \(K^n\) times the
product of deterministic comparison means \(d_v\), where \(K\ge1\).
Then
\[
 \mathbb E\!\left[\mathbf1_{\rm success}
       \bigl(|U|^{-1}\textstyle\sum_v Z_v\bigr)^n\right]
 \le K^n\bigl(|U|^{-1}\textstyle\sum_v d_v+nfL\bigr)^n.
\]
\end{lemma}

\begin{proof}
Expand the power by drawing \(n\) indices independently and uniformly
from \(U\). Mark a position as removed if its index lies in the near set
of any preceding index. On success each removed factor costs at most
\(L\); the unremoved indices satisfy the joint-comparison hypothesis.
For each removal pattern, apply that comparison, then discard the
separation conditions at the unremoved positions. Sum over the indices
in reverse order. At an unremoved position the factor is
\(|U|^{-1}\sum_v d_v\). At a removed position the necessary proximity
condition permits at most an \(nf\) fraction of indices, so its factor is
at most \(nfL\). Summing over all removal patterns gives the claimed
binomial expansion.
\end{proof}

We may first apply the joint comparison conditional on a successful
prehistory, retaining its indicator, and only afterwards discard
nonlocal restrictions for an upper bound. Typically \(nfL=o(1)\), while
\(K\) and the average of the \(d_v\)'s are bounded. Markov's inequality
at a sufficiently large constant threshold then permits a union over
\(\exp(O(n))\) columns. For \(Z_v=Np_v(y)\), the resulting bound on the
average gives
\[
 \sum_{v\in U}p_v(y)
   =\frac{|U|}{N}\left(|U|^{-1}\sum_{v\in U}Z_v\right)=o(1)
\]
when \(|U|=O(2^n)\), since \(N=C_n2^n\) and \(C_n\to\infty\).
The exceptional probability here includes the success indicator. We also
use this argument with patch-specific sizes.

\begin{lemma}[Gated posterior comparison]\label{lem:gated-posterior}
Fix any earlier information in a reference experiment. Let a candidate
\(z\) have prior law \(\pi\), and let \(F_z(t)\) be the likelihood of data
\(t\) together with a required event. Thus \(F_z\) includes the event
indicator and may have total mass less than one. Set
\[
 m(t)=\int F_z(t)\,d\pi(z).
\]
Use one common reference measure for all data densities; in a finite
space these may simply be masses. If \(Q\) is any reference probability
law on the data fixed by the earlier information and \(\epsilon>0\), then the joint
probability of the required event and
\(\{m(t)<\epsilon Q(t)\}\cup\{m(t)=0\}\) is at most \(\epsilon\).
If also \(F_z(t)\le e^sQ(t)\), the posterior
\(F_z(t)\,d\pi(z)/m(t)\) outside this exception is dominated by
\(e^s\epsilon^{-1}\pi\). Integrating the posterior against the required-event
data measure recovers the prior experiment with its gate retained.
\end{lemma}

\begin{proof}
The data measure with the required event has density \(m\). Its mass on
\(\{m<\epsilon Q\}\) is at most \(\epsilon\int Q=\epsilon\), and its
mass where \(m=0\) is zero. On the complement,
\(F_z/m\le e^s/\epsilon\). Finally, for any nonnegative test \(h\),
\[
 \int m(t)\int h(z,t)\frac{F_z(t)}{m(t)}\,d\pi(z)\,dt
   =\int\!\int h(z,t)F_z(t)\,dt\,d\pi(z),
\]
with the left integrand set to zero where \(m=0\). This is the asserted
cancellation.
\end{proof}

The gate and every selection rule in \(F_z\) must be recomputed as \(z\)
varies. To use this estimate after an assignment, we first compare the
joint law of the required outputs with the reference experiment. Only
then may we remove nonlocal prehistory restrictions, keeping the local
gate inside \(F_z\).

\subsection{Geometric height device}
\label{sec:geometric-height-device}

We next choose centers in Hamming balls while keeping the choices at
nearby sites geometrically compatible. A center has an ID consisting of
its location and level. The construction gives neighboring sites levels
that differ by at most one, and bounds the number of active IDs they can
consult. Its path-maximum rule is related to the admissible-path
construction of Lipschitz percolation~\cite[Sections~4--5]{dirr-dondl-grimmett-holroyd-scheutzow2010}.
Here eligibility may depend on all the sampled center positions, so the
independence needed in the proof comes only after separating those
positions from the later activations.

Let the queried sites form a subset \(S\) of \(Q_d\), where
\(d=\Theta(n)\). A horizontal step may join queried sites at Hamming
distance at most a fixed positive integer \(D\), including a step that
stays at the same site. Prospective center locations range over all of
\(Q_d\). Write \(V\) for the volume of a radius-\(r\) ball in \(Q_d\).
For levels \(0,\ldots,H\), the random experiment has three stages:
\begin{enumerate}[label=(\roman*),leftmargin=*]
\item At each vertex of \(Q_d\) and each level, independently, put a
prospective center with probability \(\lambda/V\), where
\(\lambda=n^{J_0}\) and the fixed constant \(J_0>2\). Write \(P\) for
all these prospective positions.
\item Using \(P\) and any auxiliary information, specify for every
site-level \((v,j)\) an eligible set \(E(v,j)\) of prospective centers in
its same-level \(r\)-ball. The required size is \(|E(v,j)|\ge\lambda/3\).
\item After those choices, activate each prospective center independently
with probability \(n^{b_0}/\lambda\). Write \(A\) for these activation
variables, where \(0<b_0<b<1\).
\end{enumerate}
A site-level is \emph{bad} if its eligible set contains no active center,
or if its same-level \((r+D)\)-ball contains more than \(n^b\) active
centers. We allow either of the following radii:
\begin{itemize}
\item \(r=\Theta(d)\) with \(r\le(1/4+o(1))d\);
\item \(r=\lfloor n^{1-\rho}\rfloor\), for fixed \(0<\rho<1\), with
\(b>b_0+(D+1)\rho\).
\end{itemize}
Choose fixed \(0<\sigma<\zeta<1\), \(0<\theta<1\), and \(a>0\) such that
\[
 b_0>a>\zeta+\sigma+(1-\theta),\qquad b>a+4\sigma.
\]

A space-level path stays in levels \([0,H]\). Its permitted steps are an
upward vertical step from a bad site-level, or a downward step of one
level together with a horizontal step. For a queried site \(v\), restrict
the whole path to spatial distance \(2DH\) from \(v\). Define the
\emph{long height} \(h(v)\) to be the largest \(j\) for which \((v,j)\)
is reachable by such a path from any level-zero start. A rule with a
shorter consultation radius restricts the same paths to the smaller ball.
If the resulting height equals \(H\) or its site-level is bad, the rule
makes no selection.

\begin{lemma}[Local height selection]\label{lem:height-selection}
In this experiment, take
\[
 R_0=\lceil\log^2 n\rceil,\qquad
 R_i=MR_{i-1},\qquad M=\lceil n^\sigma\rceil,
\]
and let \(H=R_h\) be the first scale at least \(n^{1-\zeta}\). Thus
\(H=O(n^{1-\zeta+\sigma})\).
\begin{enumerate}[label=(\roman*),leftmargin=*]
\item Except on an event of probability at most \(\exp(-n^{1+c})\), for
some fixed \(c>0\), whenever all eligible sets have the required sizes,
the long heights lie in \(\{0,\ldots,H-1\}\), their site-levels are good,
and their values differ by at most one on horizontal steps. Computing a
long height uses only bad statuses within spatial distance \(2DH\).
\item At a fixed site, on correct eligible-set sizes throughout its
consultation domain, the probability of positive long height is at most
\(\exp(-n^c)\) for some fixed \(c>0\). This remains true if one specified
prospective center is forced present.
\item If \(\theta>.9\), \(m=n^{\alpha+o(1)}\) for fixed \(\alpha>0\), and
\(\sqrt m\le H\), replacing the consultation radius by
\(D\lfloor\sqrt m\rfloor\) changes the path maximum with probability
\(o(\exp(-2m^{1/5}))\), on the same correct sizes. This estimate also
allows one forced prospective center.
\end{enumerate}
All bounds on correct sizes concern intersection with that event, not
conditioning on it. The local bounds in (ii) and (iii) also hold when,
for each fixed \(P\), one takes the supremum of the conditional activation
probability over fixed legal eligibilities, and only then averages over
\(P\).

At a good chosen height, independent pre-randomized site-level ties can
select a uniform eligible active ID. For fixed \(P,E\) of the required
sizes and a specified present ID \(c\) at level zero,
\[
 \Pr_{A,\,\mathrm{ties}}(h(v)=0\text{ and }c\text{ is selected}\mid P,E)
 \le 3/\lambda.
\]
The long and short rules may use the same ties and invalid-choice
conventions.
\end{lemma}

\begin{proof}
\noindent\textbf{Step 1: the path event to be estimated.}\quad
Use the metric
\[
 \operatorname{dist}((v,j),(v',j'))
   =\max\bigl(|j-j'|,\lceil d_H(v,v')/D\rceil\bigr).
\]
There are a bounded number of scales. Choose
\(1/4\le\eta_h<\cdots<\eta_0\le1/2\) with fixed positive gaps. A
scale-\(i\) failure from a specified point is a path stopped when it first
reaches metric distance \(R_i\), with net rise at least \(-\eta_iR_i\).
We will bound its probability by \(\exp(-n^aR_i^\theta)\), in the uniform
form specified below. The proof first finds many separated short failures inside a long one.
Their center balls may overlap. Removing the shared regions, while
allowing a small loss in the size and crowd thresholds, makes the short
failures depend on disjoint random inputs.

Fix a deterministic domain \(\mathcal D\) of allowed site-level paths and
a deterministic set \(\mathcal C\) of location-levels that may provide
centers. Both eligibility and crowd counts use only \(\mathcal C\).
At scale \(i\), require eligibility
size at least \(s_i\lambda\) and call a crowd bad above \(t_i n^b\), where
\(s_i\) increase strictly within \([1/8,1/4]\) and \(t_i\) increase
strictly within \([1/3,3/4]\), with fixed positive gaps. Let
\(F_i(P,E,A)\) be the scale failure from the specified start in these
domains with these thresholds. With \(\mathcal D,\mathcal C\) fixed, define
\[
 B_i(P)=\sup_E\Pr_A(F_i(P,E,A)\mid P,E),
\]
where eligible sets must have the stated sizes within metric distance
\(R_i\) of the start; the supremum is zero if no such sets exist. This
defines a random variable depending only on the prospective positions.
The induction claim, uniform in the two deterministic domains, is
\begin{equation}
 \mathbb E_P B_i(P)\le\exp(-n^aR_i^\theta).
 \label{eq:source-H}
\end{equation}
The actual thresholds are stronger:
actual legal sets have size at least \(\lambda/3\), and every actual bad
site-level remains bad when the crowd threshold is lowered to
\(t_i n^b\). Hence the same upper bound holds for failures in the original
experiment on correct sizes.

\noindent\textbf{Step 2: the smallest scale.}\quad
A scale-zero failure must contain an up-step. Without one, reaching
distance \(R_0\) requires at least \(R_0\) downward steps and gives net
rise at most \(-R_0\). An up-step identifies a bad site-level in the
metric ball of the start. There are at most
\(\exp(O(R_0\log n))\) possible such points.

For any fixed legal eligible set, a hole has conditional probability at
most \(\exp(-s_0n^{b_0})\). The crowd test does not depend on eligibility.
Unconditionally its count is binomial, with mean at most
\(n^{b_0}O(1)\) in the linear-radius case and
\(n^{b_0}O(n^{D\rho})\) in the sublinear case, by successive Hamming-layer
ratios. The assumed gap below \(b\) gives a binomial upper tail with
polynomial slack. The same union bound is valid after taking the
eligibility supremum: the hole bound is uniform in the eligible set,
and the crowd bound is independent of it. Since \(b_0>a\), these tails
and the count of possible points prove the induction at scale zero.

\noindent\textbf{Step 3: separated child paths and ball intersections.}\quad
Put \(R'=R_{i-1}\). A scale-\(i\) failure contains at least
\(q=\lfloor c'M\rfloor\) starts of scale-\((i-1)\) failures separated
by \(KR'\), with all child paths still inside the parent's metric ball.
Here the child paths temporarily use the parent's bad statuses; \(K\)
is a large fixed constant and \(c'>0\) is sufficiently small.

To see this, take a maximal separated set of child-failure starts. If it
has fewer than \(q\) members, their \(KR'\)-balls cover every such start.
Follow the parent path. Whenever it enters the cover, skip to its last
visit to the currently hit ball, accounting separately for one outgoing
step if necessary. No ball is used twice. These skips are bookkeeping:
each has endpoint displacement, and hence possible upward displacement,
at most \(2KR'\). Outside the cover, partition the path into pieces
stopped on first reaching distance \(R'\), with one unfinished piece per
intervening stretch. Every full piece has rise less than
\(-\eta_{i-1}R'\), since its start is not a child-failure start. The parent
displacement forces at least \(M-O((q+1)K)\) full pieces. Its total rise
is consequently at most
\[
 -\eta_{i-1}MR'+O((q+1)KR'),
\]
which is less than \(-\eta_iMR'\) when \(c'\) is small relative to the
gap \(\eta_{i-1}-\eta_i\). This contradicts the parent failure.

There are at most \(\exp(O(qR_i\log n))\) choices of these starts.
A child's center domain has levels within \(R'\) of its start and spatial
radius at most \(r+DR'+D\). For two separated children, either their
level intervals are disjoint or their spatial centers have distance
\(s\) at least a large multiple of \(R'\). At each common level, the
intersection of their enlarged balls has volume
\(V\exp(-\Omega(R'))\).

For this last estimate, a uniform point of one enlarged ball lies within
\(s/10\) of its outer layer except with probability
\(\exp(-\Omega(s))\), by the binomial-layer ratios. If it also lies in
the other ball, it must then flip at least \(.44s\) of the \(s\)
coordinates on which the centers differ. The enlarged radius is at most
\((1/4+o(1))d\), so this is a hypergeometric upper tail of probability
\(\exp(-\Omega(s))\). For sublinear radius, the enlarged ball has
radius at most \(n^{1-c''}\) for some fixed \(c''>0\); layer ratios and
a subset union bound for the differing coordinates improve both
exponents to \(\Omega(s\log n)\). Enlarging the original ball costs a
factor \(\exp(O(R'))\), or \(\exp(O(R'\log n))\) in the sublinear case.
Taking \(K\) large makes the intersection bound valid relative to \(V\)
in either case.

\noindent\textbf{Step 4: removing shared centers.}\quad
For a fixed configuration of child starts, delete every pairwise overlap
from each child center domain. The remaining regions are deterministic
and disjoint. Choose a small fixed \(\varepsilon>0\) below all the gaps
\(s_i-s_{i-1}\) and \(t_i-t_{i-1}\). Except on an event of probability
\(\exp(-\Omega(n^bR'/q))\) per pair and level, an overlap contains at most
\(\varepsilon\lambda/q\) prospective centers and
\(\varepsilon n^b/q\) active centers. Indeed their means are at most
\(\lambda e^{-\Omega(R')}\) and \(n^{b_0}e^{-\Omega(R')}\), respectively,
so the claimed estimates follow from binomial upper tails.

If none of these count exceptions occurs, each child loses at most
\(\varepsilon\lambda\) eligible centers and \(\varepsilon n^b\) active
centers. Its eligible sets still have size at least
\(s_{i-1}\lambda\). A hole stays a hole, while a crowd above \(t_i n^b\)
remains above \(t_{i-1}n^b\). Thus all the child failures transfer to the
private regions with the child thresholds.

The probability calculation must not condition on successful overlap
counts. Let \(C\) be the event that the chosen child configuration occurs,
let \(B_{\rm pos},B_{\rm act}\) be the prospective and active overlap
exceptions, and let \(P_\ell\) be the positions in the private region of
child \(\ell\). For every fixed legal eligibility assignment \(E\),
\[
 \begin{aligned}
 \Pr_A(C\mid P,E)
 &\le\mathbf1_{B_{\rm pos}}+\Pr_A(B_{\rm act}\mid P)
          +\prod_\ell b_\ell(P_\ell),\\
 b_\ell(P_\ell)
 &=\sup_{E_\ell}\Pr_A(\text{private child failure}\mid P_\ell,E_\ell).
 \end{aligned}
\]
The transferred events depend on disjoint activations once \(P,E\) are
fixed. Their probabilities are bounded by the displayed private
suprema, even if \(E\) used positions outside those regions. After
paying for the overlap exceptions we have discarded their complements,
so no conditioning remains to destroy independence. The inequality is
uniform in \(E\); taking its supremum and averaging \(P\) factors the
last product, because the \(P_\ell\)'s are independent. Induction bounds
the product of their expectations by \(\exp(-q n^a(R')^\theta)\).

The exponent exceeds both the logarithm of the configuration count and
the desired parent exponent. In fact, using \(R_i=MR'\), \(R'\le n\),
and \(q\asymp M\),
\[
 \frac{n^a(R')^\theta}{R_i\log n}
 \ge\frac{n^{a-\sigma-(1-\theta)+o(1)}}{\log n}\longrightarrow\infty,
 \qquad \frac{q}{M^\theta}\longrightarrow\infty.
\]
The ratio of an overlap-error exponent to the child exponent is at least
\(\Omega(n^{b-a-2\sigma}(R')^{1-\theta})\), which also tends to infinity
because \(b>a+4\sigma\). The number of pair-level exceptions is
polynomial. The union over configurations therefore proves
\eqref{eq:source-H} and the position-averaged induction claim.

If one prospective center is forced present, erase its location-level
from the deterministic center domain. The remaining positions have their
original independent laws. Each eligible set and crowd loses at most one
center, and
\(\lambda/3-1\ge s_i\lambda\), \(n^b-1\ge t_i n^b\) for large \(n\).
Holes persist. Every original bad-path witness therefore remains a
failure at the relaxed thresholds, uniformly in the erased center's
activation. The same induction applies.

\noindent\textbf{Step 5: heights and local approximations.}\quad
At the top scale,
\(n^aH^\theta\ge n^{a+(1-\zeta)\theta}\), whose exponent is greater
than one. Thus \eqref{eq:source-H} permits a union over all sites with
error \(\exp(-n^{1+c})\). A path from level zero that first displaces by
\(H\) would be a scale-\(h\) failure, since its endpoint still has
nonnegative level. Outside the exceptional event no such path exists.
The spatially unrestricted path maximum is therefore below \(H\), and
all paths relevant to a queried maximum lie in its \(2DH\) consultation
tube. A maximizing site-level is good, since otherwise the path could
extend upward. If neighboring maxima differed by more than one, a
downward horizontal step from the higher one would contradict maximality
at the lower site. This proves (i).

For a fixed query, classify paths reaching it by their maximum metric
displacement from their level-zero start. Below \(R_0\), positive height
requires an up-step at one of \(\exp(O(R_0\log n))\) possible points,
so Step 2 applies. For displacements in \([R_i,R_{i+1})\), the first
prefix reaching distance \(R_i\) is a scale-\(i\) failure. There are at
most \(\exp(O(R_{i+1}\log n))\) possible starts near the query. Hence
these paths have total probability at most
\[
 \sum_{i<h}\exp\bigl(O(R_{i+1}\log n)-n^aR_i^\theta\bigr).
\]
Displacements at least \(H\) contribute at most
\(\exp(O(n)-n^aH^\theta)\). The exponent comparisons in Step 4 show
that the sum, together with the base-scale contribution, is at most
\(\exp(-n^c)\). All arguments hold in deterministic restricted path
domains and for position-averaged eligibility suprema. This proves (ii)
with the stated local and forced-center uniformity.

A path reaching the query but leaving its
\(D\lfloor\sqrt m\rfloor\)-tube has displacement from its start at
least \(\sqrt m/3\). Only scales with
\(R_i\ge\sqrt m/(3M)\), or a top-scale failure, need be counted. The
previous sum now has exponent much larger than \(m^{1/5}\), since
\[
 \frac{n^a(\sqrt m/M)^\theta}{m^{1/5}}
  =n^{a-\sigma\theta+\alpha(\theta/2-1/5)+o(1)}\longrightarrow\infty.
\]
Here \(a>\sigma+(1-\theta)>\sigma\theta\) and \(\theta>.9\).
This proves (iii).

Finally, ignore the height and crowd requirements when bounding the
probability of selecting a specified level-zero ID. For fixed \(P,E\),
independent activation and a uniform tie give each member of \(E(v,0)\)
the same probability of selection, at most \(1/|E(v,0)|\le3/\lambda\).
Reinstating the requirements only decreases the joint probability.
Couple the long and short constructions with the same positions,
activations, eligibilities and ties. Whenever their path maxima agree,
the selected ID and its validity agree as well, as required when a short
rule is used inside a hypothetical-candidate calculation.
\end{proof}

\subsection{Near-product injections}
\label{sec:near-product-injections}

The two samplers below use finite row sets and finite output spaces. Row
\(i\) has a full output \(o\in\Omega_i\), with a specified label
coordinate in a common finite label set. Write \(\widehat p_i(o)\) for
its full-output law and \(p_i(y)\) for the label marginal. When there is
no side data these are the same law. An assignment is injective when its
label coordinates are distinct. A pointwise upper bound for prescribed
outputs can be summed against any nonnegative function of those outputs;
this is the form of comparison used later.

Preserving individual marginals while choosing distinct representatives is
an established feature of bipartite dependent rounding; see Gandhi,
Khuller, Parthasarathy and Srinivasan~\cite[Properties~(P1)--(P3)]{gandhi-khuller-parthasarathy-srinivasan2006}.
Their negative-correlation guarantee concerns edges incident to a common
vertex. The next lemma additionally controls prescribed outputs on
arbitrary distinct rows, with a relative error even when their
probabilities are very small.

\begin{lemma}[Calibrated near-product injections]\label{lem:near-product-injection}
Let a finite set of rows have laws on \(d\) labels, with
\(p_i(y)\le d^{-.95}\) and \(\sum_i p_i(y)\le .4\) for every label
\(y\). For sufficiently large \(d\), there is an injective random
assignment \((O_i)_i\) with exact individual full-output marginals
\(\widehat p_i\) and, for every set \(J\) of \(l\le d^{.025}\) queried
rows and prescribed outputs \((o_i)_{i\in J}\),
\begin{equation}
 \Pr(O_i=o_i\text{ for all }i\in J)
 \le \exp(d^{-.04}l)\prod_{i\in J}\widehat p_i(o_i).
 \label{eq:source-17}
\end{equation}
Targets with repeated labels have probability zero. After sampling the
labels, side data can be drawn independently from their conditional laws
given those labels. We use this sampler for assignments within small bins.
\end{lemma}

\begin{proof}
It suffices to construct the labels: the independent conditional draws
then give both assertions for full outputs.

\noindent\textbf{Step 1: complete the columns and order the rows.}\quad
First allow label atoms \(O(d^{-.95})\) and column sums at most \(1/2\).
If there are \(m\) real rows, then \(m\le d/2\). Put
\(t=\lceil2d/3\rceil\) and write \(c_y=\sum_{i=1}^m p_i(y)\).
Add \(t-m\) identical dummy rows with law
\[
 p_{\rm dum}(y)=\frac{t/d-c_y}{t-m}.
\]
The numerator is nonnegative and sums to \(t-m\), so this is a probability
law. Its atoms are \(O(1/d)\), and every completed column sums to \(t/d\).

Choose an order of the \(t\) rows, then number them in that order, such
that every column prefix satisfies
\[
 \sum_{k<j}p_k(y)=s_j+O(d^{-1/8}),\qquad s_j=(j-1)/d.
\]
A uniform random order has this property with positive probability.
Indeed the expectation is \(s_j\), and exposure without replacement has
martingale increments bounded by the column atom range. Bounded-difference
concentration~\cite{hoeffding1963,azuma1967}, followed by a union over
\(j,y\), gives the asserted error. Fix one such order.

\noindent\textbf{Step 2: track the mass already used.}\quad
At step \(j\), sample row \(j\)'s law conditional on its label still
being free. For each row \(a\), let
\[
 D_j^a=\sum_{k<j}p_a(Y_k),\qquad
 e_j=\max_{a,\,k\le j}|D_k^a-s_k|.
\]
Thus \(D_j^a\) is the mass that row \(a\) assigns to labels used before
step \(j\). Stop at the first \(e_j>1/20\); any undefined continuation
after a stop may be declared a failure. We claim that, with probability
\(1-\exp(-d^{\Omega(1)})\), there is no stop and
\(D_j^a=s_j\pm d^{-.1}\) for all \(a,j\).

Before stopping, all sampling denominators are bounded away from zero.
Conditional on the first \(j-1\) draws, the drift of \(D_j^a\) is
\[
 \mathbb E[D_{j+1}^a-D_j^a\mid Y_1,\ldots,Y_{j-1}]
  =\sum_{y\ \mathrm{free}}\frac{p_a(y)p_j(y)}{1-D_j^j}.
\]
Its martingale increments are \(O(d^{-.95})\). Bounded-increment
concentration gives cumulative errors \(O(d^{-1/8})\), simultaneously
for all rows and times, outside an event of probability
\(\exp(-d^{\Omega(1)})\).

When summing the drifts through step \(b-1\), replacing \(1-D_j^j\) by
\(1-s_j\) costs at most
\[
 O\left(\sum_y p_a(y)\sum_{j<b}e_jp_j(y)\right)
   =O\left(d^{-1/8}+d^{-1}\sum_{j<b}e_j\right).
\]
For the equality, use summation by parts on the column-prefix error
against the bounded nondecreasing sequence \(e_j\). A second summation
by parts replaces \(p_j(y)\) by \(1/d\) in the remaining drift sum:
the sequence \(\mathbf1_{\{y\ \mathrm{free\ before}\ j\}}/(1-s_j)\)
has bounded variation. The resulting sum is
\[
 d^{-1}\sum_{j<b}\frac{1-D_j^a}{1-s_j}+O(d^{-1/8}).
\]
Subtracting \(s_b\) now gives, even at a potential stopping index,
\[
 e_b\le Kd^{-1/8}+\frac Kd\sum_{j<b}e_j.
\]
Iteration bounds \(e_b\) by \(O(d^{-1/8})\). This is smaller than both
\(1/20\) and \(d^{-.1}\) for large \(d\), proving the claim. Call this
successful tracking event \(G\), and condition the sampler on \(G\).

\noindent\textbf{Step 3: compare specified targets multiplicatively.}\quad
Fix \(l\) distinct positive-probability target labels \(y_i\) on queried
rows. Define a second sampler that forces \(Y_i=y_i\) at each queried
step and, at every ordinary step, excludes all target labels whose
queried steps are still pending. On a path realizing the targets and
satisfying \(G\), the likelihood of the original unconditioned sampler
divided by that of the forcing sampler is the product of
\[
 \frac{p_i(y_i)}{1-D_i^i}\quad(i\text{ queried}),\qquad
 1-\frac{\sum_{i>j,\ i\ \mathrm{queried}}p_j(y_i)}{1-D_j^j}
       \quad(j\text{ ordinary}).
\]
After extracting \(\prod_i p_i(y_i)\), the target-step denominators and
the ordinary-step factors cancel to small relative error. For each target,
the prefix estimate and summation by parts give
\[
 \sum_{j<i}\frac{p_j(y_i)}{1-s_j}
       =-\log(1-s_i)+O(d^{-1/8}).
\]
Tracking permits replacement of \(s_j\) by \(D_j^j\) at error
\(O(d^{-.1})\); omitting queried steps costs \(O(l d^{-.95})\) per
target. Thus the linear terms in the ordinary-step logarithms cancel
\(-\log(1-D_i^i)\). The quadratic terms have total size
\(O(l^2d^{-.95})\), because each subtracted fraction is
\(O(l d^{-.95})\) and their sum is \(O(l)\). The total log-error per
target is \(O(d^{-.1}+l d^{-.95})\).

Integrating this ratio under the forcing law, and dividing by
\(\Pr(G)=1-\exp(-d^{\Omega(1)})\), gives the joint upper comparison
\(\exp(O(d^{-.09}l))\prod_i p_i(y_i)\). For a single target the same
ratio also gives a lower comparison \((1-O(d^{-.09}))p_i(y_i)\): the
forcing sampler itself tracks with exponentially small failure
probability. Indeed, before a constant-error stop, excluding the one
pending label changes a step's distribution in total variation by
\(O(p_j(y_i))\). Since the tracked increment is \(O(d^{-.95})\), the
cumulative drift change is
\(O(d^{-.95}\sum_jp_j(y_i))=O(d^{-.95})\); the forced step contributes
another \(O(d^{-.95})\). The proof of Step 2 therefore still applies.
The target mass was extracted before this tracking exception was bounded;
no absolute error is divided by a tiny atom.

\noindent\textbf{Step 4: calibrate the individual marginals exactly.}\quad
Return to the original \(.4\) column bound. For arbitrary real prices
\(c_i(y)\) on the real row-label pairs, put
\(\bar c_i=\sum_y p_i(y)c_i(y)\) and \(\delta=d^{-.05}\). Perturb row
\(i\) by multiplying its probabilities by
\(1+\delta\operatorname{sign}(c_i(y)-\bar c_i)\), then normalize.
Writing \(p'_i\) for the result, its gain in centered expected price is
\[
 \begin{aligned}
 \sum_y p'_i(y)(c_i(y)-\bar c_i)
 &=\frac{\delta\sum_y p_i(y)|c_i(y)-\bar c_i|}
        {1+\delta\sum_y p_i(y)\operatorname{sign}(c_i(y)-\bar c_i)}
 \\
 &\ge\tfrac12\delta\sum_y p_i(y)|c_i(y)-\bar c_i|.
 \end{aligned}
\]
The perturbed rows satisfy the relaxed hypotheses of Step 1. The
tracking sampler's two-sided singleton error is \(O(d^{-.09})\) relative
to \(p'_i\), so its error in this centered price is at most
\(O(d^{-.09})\sum_y p_i(y)|c_i(y)-\bar c_i|\), smaller than the gain.
Its joint upper bound, measured against the \emph{original} inputs, costs
\(\exp(O(d^{-.05}l+d^{-.09}l))\), which fits
\eqref{eq:source-17} for large \(d\).

Let \(\mathcal C\) be the set of all laws on injections of the real rows
satisfying the joint upper inequalities relative to those original
inputs. These are linear inequalities on a finite probability simplex,
so \(\mathcal C\) is compact and convex. Every price-dependent
construction just made belongs to this same set, and its marginal vector
has price at least that of the original marginal vector. If the original
vector were outside the marginal image of \(\mathcal C\), a separating
linear functional would contradict that property. It therefore belongs
to the image, giving exact marginals while preserving all joint bounds.
\end{proof}

\subsection{Clock lemma}
\label{sec:clock-lemma}

The next sampler also enforces local restrictions on full outputs. It
gives a joint upper comparison with the product law, without the exact
marginals of Lemma~\ref{lem:near-product-injection}. We retain the finite
row and full-output spaces, and the notation \(\widehat p_a,p_a\), from
the preceding subsection. A predicate's \emph{scope} is the set of rows
whose full outputs it reads.

\begin{lemma}[Clock sampling]\label{lem:clock-sampling}
Let there be \(g\) labels, with \(\log g=O(n)\), and finitely many
probability rows \(\widehat p_a\). For each fixed \(B\ge1\), there are
sufficiently large constants \(A_*,P_*\) and a small absolute constant
\(\theta_0>0\) with the following property. Suppose
\[
 \sum_a p_a(y)\le\theta_0,\qquad p_a(y)\le n^{-A_*}
\]
for every label and row. Let a finite list of failure predicates have
probability at most \(n^{-P_*}\) under independent row outputs. Assume
each predicate reads at most \(n^B\) rows and each row belongs to at most
\(n^B\) predicate scopes. Then there is a random assignment with distinct
labels avoiding all predicates, such that for every set \(J\) of at most
\(n^B\) queried rows,
\[
 \Pr(O_a=o_a\text{ for all }a\in J)
       \le(1+o(1))\prod_{a\in J}\widehat p_a(o_a).
\]
The error is uniform over the rows and targets for fixed \(B\), the chosen
constants, and a fixed bound in \(\log g=O(n)\). Noninjective targets
have probability zero. Unneeded predicates and an empty assignment may
be ignored.
\end{lemma}

\begin{proof}
Choose \(K_0\) sufficiently large in terms of \(B\), then choose
\(A_*,P_*\) sufficiently large in terms of \(K_0,B\). We first construct
a greedy matching from independent edge clocks and use finite backward
explorations to detect failures. The proof gives two estimates: a joint
product bound for prescribed outputs, and a small total probability for
bad exploration outcomes involving any fixed row or label. The first
controls the raw matching law. The second lets
Lemma~\ref{lem:conditional-avoidance} enforce a match at every real row
and avoidance of every predicate, at a multiplicative cost \(1+o(1)\)
on the queried outputs.

\noindent\textbf{Step 1: trim rare outputs and complete the columns.}\quad
For each row \(a\) and incident failure predicate \(F\), discard outputs
\(o\) for which
\[
 \Pr(F\mid O_a=o)>n^{-P_*/2}
\]
under the product law. Markov's inequality bounds the discarded mass
for each predicate by \(n^{-P_*/2}\), so the total loss per row is at
most \(\varepsilon=n^{B-P_*/2}\). Renormalize the retained laws. A test
or query involves at most \(n^B\) rows, so this changes its probability
by at most \((1-\varepsilon)^{-n^B}=1+o(1)\). Under the new laws, every
predicate has probability at most \(n^{-P_*/3}\), both unconditionally
and after any one row in its scope is pinned to a retained output. We
continue to denote these laws by \(\widehat p_a,p_a\).

Put \(\theta=\lceil10^{-6}g\rceil/g\), and take
\(\theta_0=10^{-8}\), for example. If there are \(m\) real rows and
\(c_y=\sum_a p_a(y)\), add \(\theta g-m\) identical dummy rows with
law \((\theta-c_y)/(\theta g-m)\). This is a probability law because
its numerator sums to the denominator. The denominator is a positive
integer of order \(g\), even after the small trimming loss. Every
completed column sums to \(\theta\), and the dummy atoms are \(O(1/g)\).
The nonempty case has \(g\ge n^{A_*}\), since an original real row is a
probability law with atoms at most \(n^{-A_*}\). Thus the largest edge
rate in the completed row-label graph is
\(\lambda_*=O(n^{-A_*})\).

\noindent\textbf{Step 2: matching and its backward explorations.}\quad
Run an independent Poisson process of rate \(p_a(y)\) on every positive
row-label edge \((a,y)\). Each arrival has an independent mark drawn
from the conditional law of the full output given label \(y\). Through
time \(T_0=K_0\log n\), greedily match an arriving edge if both endpoints
are free. Only its first arrival can matter: if that arrival is rejected,
an endpoint is already matched and never becomes free again.

There are two types of tests. A predicate test fails if all rows in its
scope are matched and their outputs violate the predicate. A real
singleton-row test fails if that row is unmatched. For each test we also
forbid an excessively large exploration of the clocks needed to determine
its outcome.

Start that exploration with the scope rows as roots, each requesting its
matching history through time \(T_0\). An endpoint's requested horizon
is the latest time up to which its status is needed. Process the pending
endpoint with largest horizon. In a fixed order, inspect its incident
edges for first arrivals by that horizon, reusing information already
revealed. A hit at time \(t\) requests the other endpoint through a time
strictly before \(t\); retain the largest request if it is requested
again. Horizons are processed in decreasing order, so a processed
endpoint can never receive a larger request. These inspections determine
the matching decisions at the roots by backward closure through all
earlier events they require.

The active endpoints are the roots and every endpoint reached by a hit,
including pending and processed endpoints. Stop with a \emph{giant}
failure as soon as their number reaches \(L=\lceil n^{.1K_0}\rceil\).
Register the endpoint of the last hit before stopping. Consequently every
inspected hit always has both endpoints active, even on a giant leaf.

\noindent\textbf{Step 3: exploration leaves and nonneighbor conditioning.}\quad
Eventually we discretize first-arrival times with a sufficiently fine
finite mesh and a fixed tie order. A leaf of an exploration then specifies
a product rectangle in the independent edge data. A hit pins an arrival
bin and mark on an edge whose two endpoints are active. A negative
inspection requires absence up to a specified cutoff on an edge with at
least one active endpoint. Every realization in this rectangle follows
the same inspections and reaches the same leaf.

Use the individual bad leaves as forbidden events, joining two leaves
when their active endpoint sets intersect. Leaves with disjoint active
sets can share an inspected edge only if both require absence on it.
To force a positive-probability leaf, condition its edges separately. On
absence edges couple the conditional first time to be no earlier than
the old time; on hit edges fix the required data; leave other edges
unchanged. Every already occurring nonneighbor leaf still occurs. This
holds simultaneously for any collection of nonneighbors, even if they
overlap with one another. If \(A\) is avoidance of such a collection
and \(E\) is the forced leaf, the coupling gives
\(\Pr(A\mid E)\le\Pr(A)\), hence
\(\Pr(E\mid A)\le\Pr(E)\) whenever defined.

It remains to show, uniformly in an endpoint \(u\), that
\[
 \sum_{\substack{E\ \mathrm{bad\ leaf}:\ u\ \mathrm{active\ in}\ E}}
       \Pr(E)\le n^{-.6K_0}.
\]
Equivalently, sum over tests the probability that the test has a bad
outcome and its exploration reaches \(u\). We prove this in continuous
time before choosing the mesh.

\noindent\textbf{Step 4: inserted paths bound the incidence sum.}\quad
Reaching \(u\) from a test's roots requires a path of time-decreasing
arrivals from one of its root rows to \(u\). Such a certificate uses
distinct edges, since the exploration inspects first arrivals; the empty
path is allowed. For a fixed edge path, count its possible distinguished
arrival tuples. Poisson insertion expresses their expected count,
with a specified bad outcome, as the integral of the product of edge
rates over ordered times in \([0,T_0]\), multiplied by the bad-outcome
probability in an ordinary realization with those arrivals inserted.
This identity follows directly by summing over the Poisson counts and
distinguishing one of their uniform-time points on each edge. Inserted
marks may be fixed anywhere in their support for uniform estimates.

To sum certificates, reverse them from \(u\) and allow all walks. Total
outgoing rates alternate between \(1\) at a row and \(\theta\) at a
label. The integrated rate sum for a start at a row is
\[
 \sum_{j\ge0}\theta^{\lfloor j/2\rfloor}\frac{T_0^j}{j!}
 =\cosh(\sqrt\theta T_0)+\theta^{-1/2}\sinh(\sqrt\theta T_0)
 =O(e^{\sqrt\theta T_0}/\sqrt\theta).
\]
Starting at a label only decreases this upper bound. Each root row
belongs to at most \(n^B+1\) tests. Lengths above
\(K_0^2\log n\) contribute less than \(n^{-K_0^2}\) before this
incidence factor, by the factorial tail of the series.

For shorter paths, first dispose of those meeting two distinct rows of
the same test, bounding the bad-outcome probability simply by one. Choose
their two positions on the path and a test incident to the earlier row.
In the normalized reversed walk, a subsequent label-to-row transition
enters that test's scope with probability at most
\(n^B\lambda_*/\theta\). Their total contribution is therefore at most
\[
 O(e^{\sqrt\theta T_0}/\sqrt\theta)(1+\log n)^{O(1)}
       (n^B+1)n^B\lambda_*/\theta.
\]
The remaining short insertions meet at most one row of the test. We now
bound bad outcomes uniformly under each such insertion.

\noindent\textbf{Step 5: the exploration rarely becomes giant.}\quad
Include the endpoints of all inserted arrivals as extra full-horizon
roots. Every reached endpoint is connected to one of these roots by
ordinary arrivals: take the part of its reaching path after the last
inserted hit. There are at most \(n^B+O(\log n)\) roots.

During largest-horizon processing, an edge from the current endpoint to
an unprocessed endpoint is fresh: neither of its endpoints could have
inspected it earlier. Superpose the fresh arrivals and, when necessary,
independent extra arrivals to reach total rate \(1\) at row endpoints
and \(\theta\) at label endpoints. They are dominated by a two-type
branching forest. A node with horizon \(t\) has Poisson children of the
opposite type at times in \([0,t]\), each child inheriting its arrival
time as horizon. To handle repeated requests for a graph endpoint,
attach its sole processing to the node making its highest request;
its offspring have not yet been generated. Unused tree nodes can be
filled with independent subtrees. This deferred construction dominates
the active endpoint count by the forest's total progeny.

For \(\delta=n^{-.01K_0}\), let \(H_{\rm row}(t)\) and
\(H_{\rm label}(t)\) be the exponential moments of single-root progeny,
including the root. They satisfy
\[
 \begin{aligned}
 H_{\rm row}(t)&=\exp\left(\delta+
                    \int_0^t(H_{\rm label}(s)-1)\,ds\right),\\
 H_{\rm label}(t)&=\exp\left(\delta+
                    \theta\int_0^t(H_{\rm row}(s)-1)\,ds\right).
 \end{aligned}
\]
Induction on truncated generations bounds their excesses above one by
\(4\delta e^{2\sqrt\theta t}/\sqrt\theta\) and
\(4\delta e^{2\sqrt\theta t}\), respectively. Substitution proves
these bounds using \(e^x-1\le2x\) for the resulting small positive
exponents. With \(R\le n^B+O(\log n)\) roots, exponential Markov gives
\[
 \Pr(\text{giant})\le
 \exp\left(-\delta L+
      O(R\delta e^{2\sqrt\theta T_0}/\sqrt\theta)\right).
\]
Here \(\delta L\ge n^{.09K_0}\), while the second term is
\(O(n^{B-(.01-2\sqrt\theta)K_0})\). Choosing
\((.1-2\sqrt\theta)K_0>B\) makes this tail superpolynomially small,
uniformly in the fixed short insertion.

\noindent\textbf{Step 6: track the free mass in a background process.}\quad
Fix a prescription removing at most \(2n^B\) endpoints and inserting
\(O(\log n)\) arrivals between the remaining endpoints. In the resulting
background matching define, even for original endpoints that were removed,
\[
 X_a(t)=\sum_{y\ \mathrm{free\ in\ background}}p_a(y),\qquad
 R_y(t)=\sum_{a\ \mathrm{free\ in\ background}}p_a(y).
\]
These are the available label mass of a row and the available row mass
at a label. The deterministic comparison functions solve
\[
 z'=-qz,\qquad q'=-\theta zq,\qquad z(0)=q(0)=1.
\]
In particular \(q=1-\theta+\theta z\) and \(0<q,z\le1\). We claim,
with error probability at most \(\exp(-n^{A_*/3})\), that throughout
\([0,T_0]\), uniformly over original endpoints,
\[
 X_a(t)=q(t)+o(n^{-3B}),\qquad
 R_y(t)=\theta z(t)+o(n^{-3B}).
\]
The estimate is uniform for each fixed removal and insertion prescription;
we do not take a union over all such prescriptions.

A free label \(y\) is matched at ordinary rate \(R_y\), and a free row
\(a\) at rate \(X_a\). Hence the ordinary drifts of these masses are
\(-\sum_{y\ \mathrm{free}}p_a(y)R_y\) and
\(-\sum_{a\ \mathrm{free}}p_a(y)X_a\). If \(E(t)\) is the maximum of
the claimed approximation errors, each drift differs from its ODE drift
by at most \(2E(t)\). Removals cause initial error
\(O(n^B\lambda_*)\), and inserted matches cause total jumps
\(O(\log n\,\lambda_*)\).

The drift-compensated ordinary-jump martingales have jumps at most
\(\lambda_*\) and predictable variance at most \(T_0\lambda_*\).
Their running absolute maxima exceed \(n^{-A_*/4}\) with probability
at most \(\exp(-n^{A_*/2-o(1)})\) per endpoint. For completeness,
compensating \(e^{vM_t}\) costs log-rate at most \(v^2\) times the
variance rate when \(|v|\lambda_*\) is small; this follows from the
exponential series for each jump. Use both signs of
\(v=n^{-A_*/4}/(2T_0\lambda_*)\) and stop at the first crossing. This
is the bounded-jump exponential estimate used here; compare~\cite{freedman1975}.
A union over \(O(g)\) endpoints is permitted by \(\log g=O(n)\).
Finally drift iteration yields
\[
 E(t)\le
 \left[O((n^B+\log n)\lambda_*)+n^{-A_*/4}\right]e^{2T_0}
 \le n^{-A_*/4+2K_0+o(1)}.
\]
Taking \(A_*/4>2K_0+3B\), with slack, proves the claim and its stated
exception probability.

\noindent\textbf{Step 7: ordinary target matches have a product upper bound.}\quad
Fix \(k\le n^B\) full-output targets with distinct labels, also allowing
any fixed short insertion from Step 4. Count tuples of ordinary arrivals
that realize all the target matches. Insert their distinguished arrivals
at candidate times \(t_1,\ldots,t_k\). Their intensities factor out the
entire product \(\prod_a\widehat p_a(o_a)\), including all prescribed
marks, before any probability estimate is made.

If these are the target matches, none of their endpoints ever matches
to an outside endpoint. The outside process therefore agrees exactly
with the background obtained by removing the queried rows and target
labels, keeping the fixed inserted arrivals between remaining endpoints.
Given this background, an ordinary arrival from a target endpoint to a
free outside endpoint is forbidden before its candidate matching time.
These are disjoint clocks independent of the background. Ignoring other
necessary conditions gives the following upper bound for the remaining
time integrand:
\[
 \mathbb E\exp\left(-\sum_{(a,y_a)}
           \int_0^{t_a}(X_a(s)+R_{y_a}(s))\,ds\right).
\]
The ODE identity
\((\log(zq))'=-(q+\theta z)\) and Step 6 bound this, on the tracking
event, by \((1+o(1))\prod_a z(t_a)q(t_a)\): its logarithmic error is
at most \(kT_0\,o(n^{-3B})=o(1)\). After integrating the times, the
ordinary-target probability is at most
\[
 \left(\prod_a\widehat p_a(o_a)\right)
 \left[(1+o(1))\left(\int_0^{T_0}z(t)q(t)\,dt\right)^k
                 +T_0^k e^{-n^{A_*/3}}\right]
 \le(1+o(1))\prod_a\widehat p_a(o_a).
\]
Indeed \(\int_0^{T_0}zq=1-z(T_0)\le1\), and the second term is
\(o(1)\) for our choice of \(A_*\). In particular the tracking exception
still carries the full target product; the bound is relative even for
arbitrarily small positive atoms.

For an unmatched singleton, remove just its row. The same absence
calculation gives \(\mathbb E\exp(-\int_0^{T_0}X_a(s)\,ds)\) as an
upper bound. Since \(q\ge1-\theta\), its probability under the short
insertion is at most \(n^{-(1-\theta)K_0+o(1)}\).

\noindent\textbf{Step 8: bound bad leaves and impose avoidance.}\quad
For the remaining short path and predicate pairs, at most one row of the
predicate lies on the path. If all its matches are ordinary, sum Step 7
over its failing output tuples. Otherwise one match uses an inserted
arrival at that row. Pin its output, discard its match requirement, and
sum the ordinary-target bounds for the other rows using the one-pinned
failure estimate from Step 1. There are \(O(\log n)\) inserted arrivals
to consider. In either case the matched-failure probability is at most
\(O(1+\log n)n^{-P_*/3}\).

Multiply these estimates and the giant bound by the path and test factor
\(O((n^B+1)n^{\sqrt\theta K_0}/\sqrt\theta)\) from Step 4. The
unmatched contribution is at most
\(n^{B-(1-\theta-\sqrt\theta)K_0+o(1)}\). The two-row exception is
at most \(n^{2B+\sqrt\theta K_0-A_*+o(1)}\), and long paths cost
\(O(n^{B-K_0^2})\). Choose \(K_0\) so that
\(B<(.4-\theta-\sqrt\theta)K_0\), in addition to Step 5's condition;
then take \(A_*,P_*\) large enough for the other terms and tracking.
Every contribution has a fixed power of slack below \(n^{-.6K_0}\).
This proves the active-endpoint bound from Step 3.

Choose the time mesh sufficiently fine for this bound to persist.
Continuous arrival times are distinct almost surely; matching decisions
and truncated explorations therefore converge almost surely as the
mesh is refined. There are finitely many tests, endpoints and full-output
target products. We may also preserve the raw joint target bound, with
no path insertions, with negligible relative loss at every positive
product. The discretized exploration now has finitely many leaves.

Assign each bad leaf charge twice its raw probability. A leaf has at
most \(L\) active endpoints, so its total neighboring charge is at most
\(2L n^{-.6K_0}=O(n^{-.5K_0})\). Thus its probability is at most its
charge times the product of its neighbors' avoidance factors, and
Lemma~\ref{lem:conditional-avoidance} supplies simultaneous avoidance.
All real rows are then matched, all listed predicates are avoided, and
all the required explorations are nongiant.

Finally partition any joint target event by the non-giant singleton
leaves of its queried rows. Their intersections are product rectangles
with at most \(n^B L\) active endpoints. They have the same forcing
property against bad leaves with disjoint active sets. The total
neighboring charge is at most \(2n^B L n^{-.6K_0}\), so the avoidance
comparison costs
\(\exp(O(n^{B-.5K_0}))=1+o(1)\). Summing the rectangles gives the raw
joint target bound from Step 7. The initial trimming normalization costs
another \(1+o(1)\), proving the assertion for the original full-output
laws.
\end{proof}
